\chapter{Release 1 --- Conception et développement}

\section*{Introduction}
\addcontentsline{toc}{section}{Introduction}
Cette release constitue le c\oe{}ur métier de la plateforme de facturation. Elle couvre le développement de l'ensemble des microservices fonctionnels répartis sur trois sprints. Chaque sprint est accompagné de sa modélisation UML (diagrammes de cas d'utilisation, de classes et de séquence).

%======================================================================
% SPRINT 1
%======================================================================
\section{Sprint 1 --- Boutique-service et Customer-service}

\subsection{Objectif du sprint}
Ce sprint vise à mettre en place les deux premiers microservices fondamentaux : la gestion des boutiques (points de vente) et la gestion des clients.

\subsection{Sprint Backlog}

\begin{table}[H]
\centering
\renewcommand{\arraystretch}{1.2}
\begin{tabular}{|c|p{9cm}|c|c|}
\hline
\textbf{ID} & \textbf{Tâche} & \textbf{Priorité} & \textbf{Points} \\
\hline
1.1 & CRUD Boutique (création, consultation, modification) & Must & 5 \\
\hline
1.2 & CRUD Client avec validation CIN (8 chiffres) et téléphone (8 chiffres) & Must & 8 \\
\hline
1.3 & Modification et suspension d'un client & Must & 3 \\
\hline
1.4 & Publication d'événements client via le bus de messages & Should & 3 \\
\hline
\end{tabular}
\caption{Sprint Backlog --- Sprint 1}
\label{tab:sprint1_backlog}
\end{table}

\subsection{Diagramme de cas d'utilisation --- Sprint 1}

\begin{figure}[H]
    \centering
    \fbox{\parbox{0.85\textwidth}{\centering\vspace{2.5cm}\textbf{[Diagramme de cas d'utilisation --- Sprint 1]}\\\vspace{0.2cm}\textit{Gestion des boutiques et gestion des clients}\vspace{2.5cm}}}
    \caption{Diagramme de cas d'utilisation du Sprint 1}
    \label{fig:uc_sprint1}
\end{figure}

\subsection{Diagramme de classes --- Sprint 1}

\subsubsection{Boutique-service}
Le microservice Boutique gère les points de vente et le stock de cartes SIM. Les principales entités sont :
\begin{itemize}
    \item \textbf{Boutique :} Représente un point de vente avec un nom, une adresse, un responsable et un statut (ACTIVE, INACTIVE).
    \item \textbf{SimCard :} Représente une carte SIM avec un numéro ICCID, un statut (AVAILABLE, ASSIGNED, SUSPENDED) et une référence au client assigné.
\end{itemize}

\begin{figure}[H]
    \centering
    \fbox{\parbox{0.85\textwidth}{\centering\vspace{2.5cm}\textbf{[Diagramme de classes --- Boutique-service]}\\\textit{Boutique, SimCard}\vspace{2.5cm}}}
    \caption{Diagramme de classes du Boutique-service}
    \label{fig:class_boutique}
\end{figure}

\subsubsection{Customer-service}
Le microservice Customer gère le cycle de vie des clients. L'entité principale est :
\begin{itemize}
    \item \textbf{Customer :} Contient les informations personnelles du client (nom, prénom, email, CIN, téléphone, adresse), une référence unique (format \texttt{CLT-ANNÉE-UUID}), un statut (ACTIVE, SUSPENDED, TERMINATED) et la référence de la boutique rattachée.
\end{itemize}

\begin{figure}[H]
    \centering
    \fbox{\parbox{0.85\textwidth}{\centering\vspace{2.5cm}\textbf{[Diagramme de classes --- Customer-service]}\\\textit{Customer}\vspace{2.5cm}}}
    \caption{Diagramme de classes du Customer-service}
    \label{fig:class_customer}
\end{figure}

\subsection{Diagramme de séquence --- Créer un client}

\begin{figure}[H]
    \centering
    \fbox{\parbox{0.85\textwidth}{\centering\vspace{2.5cm}\textbf{[Diagramme de séquence --- Créer un client]}\\\textit{Agent commercial $\rightarrow$ API Gateway $\rightarrow$ Customer-service $\rightarrow$ Base de données}\vspace{2.5cm}}}
    \caption{Diagramme de séquence : création d'un client}
    \label{fig:seq_create_customer}
\end{figure}

\subsection{Interfaces réalisées --- Sprint 1}
\begin{figure}[H]
    \centering
    \fbox{\parbox{0.85\textwidth}{\centering\vspace{2.5cm}\textbf{[Capture d'écran --- Liste des boutiques]}\vspace{2.5cm}}}
    \caption{Interface de gestion des boutiques}
    \label{fig:ui_boutiques}
\end{figure}

\begin{figure}[H]
    \centering
    \fbox{\parbox{0.85\textwidth}{\centering\vspace{2.5cm}\textbf{[Capture d'écran --- Liste des clients]}\vspace{2.5cm}}}
    \caption{Interface de gestion des clients}
    \label{fig:ui_clients}
\end{figure}

%======================================================================
% SPRINT 2
%======================================================================
\section{Sprint 2 --- Catalog-service et Subscription-service}

\subsection{Objectif du sprint}
Ce sprint vise à développer le catalogue des services télécom et le mécanisme de souscription aux offres commerciales.

\subsection{Sprint Backlog}

\begin{table}[H]
\centering
\renewcommand{\arraystretch}{1.2}
\begin{tabular}{|c|p{9cm}|c|c|}
\hline
\textbf{ID} & \textbf{Tâche} & \textbf{Priorité} & \textbf{Points} \\
\hline
2.1 & CRUD Service télécom (voix, data, SMS) & Must & 5 \\
\hline
2.2 & CRUD Offre commerciale (regroupement de services) & Must & 8 \\
\hline
2.3 & Souscription d'un client à une offre (création d'abonnement) & Must & 13 \\
\hline
2.4 & Gestion du cycle de vie d'un abonnement (activation, suspension, résiliation) & Must & 5 \\
\hline
\end{tabular}
\caption{Sprint Backlog --- Sprint 2}
\label{tab:sprint2_backlog}
\end{table}

\subsection{Diagramme de cas d'utilisation --- Sprint 2}

\begin{figure}[H]
    \centering
    \fbox{\parbox{0.85\textwidth}{\centering\vspace{2.5cm}\textbf{[Diagramme de cas d'utilisation --- Sprint 2]}\\\textit{Gestion du catalogue et des abonnements}\vspace{2.5cm}}}
    \caption{Diagramme de cas d'utilisation du Sprint 2}
    \label{fig:uc_sprint2}
\end{figure}

\subsection{Diagramme de classes --- Sprint 2}

\subsubsection{Catalog-service}
Le microservice Catalog gère les services télécom et les offres commerciales :
\begin{itemize}
    \item \textbf{ServiceEntity :} Représente un service unitaire (voix, data, SMS) avec un code, un nom, une description et un tarif.
    \item \textbf{Offre :} Regroupe un ensemble de services en une offre commerciale, avec un nom, un prix, une durée de validité et un statut (ACTIVE, DISCONTINUED).
\end{itemize}

\begin{figure}[H]
    \centering
    \fbox{\parbox{0.85\textwidth}{\centering\vspace{2.5cm}\textbf{[Diagramme de classes --- Catalog-service]}\\\textit{ServiceEntity, Offre}\vspace{2.5cm}}}
    \caption{Diagramme de classes du Catalog-service}
    \label{fig:class_catalog}
\end{figure}

\subsubsection{Subscription-service}
Le microservice Subscription gère les abonnements :
\begin{itemize}
    \item \textbf{Abonnement :} Lie un client à une offre avec une date de début, une date de fin, un statut (ACTIVE, SUSPENDED, TERMINATED) et une fréquence de facturation.
\end{itemize}

\begin{figure}[H]
    \centering
    \fbox{\parbox{0.85\textwidth}{\centering\vspace{2.5cm}\textbf{[Diagramme de classes --- Subscription-service]}\\\textit{Abonnement}\vspace{2.5cm}}}
    \caption{Diagramme de classes du Subscription-service}
    \label{fig:class_subscription}
\end{figure}

\subsection{Diagramme de séquence --- Souscrire un abonnement}

\begin{figure}[H]
    \centering
    \fbox{\parbox{0.85\textwidth}{\centering\vspace{2.5cm}\textbf{[Diagramme de séquence --- Souscrire un abonnement]}\\\textit{Agent $\rightarrow$ Gateway $\rightarrow$ Subscription-service $\rightarrow$ Customer-service (Feign) $\rightarrow$ BD}\vspace{2.5cm}}}
    \caption{Diagramme de séquence : souscription d'un abonnement}
    \label{fig:seq_subscribe}
\end{figure}

\subsection{Interfaces réalisées --- Sprint 2}
\begin{figure}[H]
    \centering
    \fbox{\parbox{0.85\textwidth}{\centering\vspace{2.5cm}\textbf{[Capture d'écran --- Catalogue des offres]}\vspace{2.5cm}}}
    \caption{Interface du catalogue des offres}
    \label{fig:ui_catalog}
\end{figure}

%======================================================================
% SPRINT 3
%======================================================================
\section{Sprint 3 --- Billing-service, Auth, API Gateway et Frontend}

\subsection{Objectif du sprint}
Ce sprint finalise le développement métier avec le service de facturation, l'authentification JWT, le gateway centralisé et l'interface Angular.

\subsection{Sprint Backlog}

\begin{table}[H]
\centering
\renewcommand{\arraystretch}{1.2}
\begin{tabular}{|c|p{9cm}|c|c|}
\hline
\textbf{ID} & \textbf{Tâche} & \textbf{Priorité} & \textbf{Points} \\
\hline
3.1 & Génération de factures pour un abonnement & Must & 13 \\
\hline
3.2 & Export de facture en PDF professionnel & Must & 5 \\
\hline
3.3 & Envoi de facture par email via Brevo SMTP & Should & 5 \\
\hline
3.4 & Authentification JWT avec contrôle RBAC & Must & 5 \\
\hline
3.5 & API Gateway avec routage et filtrage JWT & Must & 8 \\
\hline
3.6 & Interface Angular (boutiques, clients, abonnements, factures) & Must & 13 \\
\hline
3.7 & Eureka Server pour la découverte des services & Should & 3 \\
\hline
\end{tabular}
\caption{Sprint Backlog --- Sprint 3}
\label{tab:sprint3_backlog}
\end{table}

\subsection{Diagramme de cas d'utilisation --- Sprint 3}

\begin{figure}[H]
    \centering
    \fbox{\parbox{0.85\textwidth}{\centering\vspace{2.5cm}\textbf{[Diagramme de cas d'utilisation --- Sprint 3]}\\\textit{Facturation, authentification et interface}\vspace{2.5cm}}}
    \caption{Diagramme de cas d'utilisation du Sprint 3}
    \label{fig:uc_sprint3}
\end{figure}

\subsection{Diagramme de classes --- Sprint 3}

\subsubsection{Billing-service}
Le microservice Billing gère le cycle de vie des factures :
\begin{itemize}
    \item \textbf{Facture :} Contient le numéro de facture, la référence client, le montant HT, le taux de TVA (19\%), le montant TTC, le statut (DRAFT, SENT, PAID, OVERDUE, CANCELLED) et les dates d'émission et d'échéance.
    \item \textbf{LigneFacture :} Représente une ligne dans la facture avec la description du service, la quantité, le prix unitaire et le montant.
\end{itemize}

\begin{figure}[H]
    \centering
    \fbox{\parbox{0.85\textwidth}{\centering\vspace{2.5cm}\textbf{[Diagramme de classes --- Billing-service]}\\\textit{Facture, LigneFacture}\vspace{2.5cm}}}
    \caption{Diagramme de classes du Billing-service}
    \label{fig:class_billing}
\end{figure}

\subsubsection{Authentication-service}
Le microservice Authentication gère l'accès sécurisé :
\begin{itemize}
    \item \textbf{User :} Contient le nom d'utilisateur, le mot de passe chiffré, le rôle (ADMIN, RESPONSABLE\_BOUTIQUE, AGENT\_COMMERCIAL) et le statut du compte.
\end{itemize}

\begin{figure}[H]
    \centering
    \fbox{\parbox{0.85\textwidth}{\centering\vspace{2.5cm}\textbf{[Diagramme de classes --- Authentication-service]}\\\textit{User, Role (enum)}\vspace{2.5cm}}}
    \caption{Diagramme de classes de l'Authentication-service}
    \label{fig:class_auth}
\end{figure}

\subsection{Diagramme de séquence --- Générer une facture}

\begin{figure}[H]
    \centering
    \fbox{\parbox{0.85\textwidth}{\centering\vspace{2.5cm}\textbf{[Diagramme de séquence --- Générer une facture]}\\\textit{Agent $\rightarrow$ Gateway $\rightarrow$ Billing-service $\rightarrow$ Subscription (Feign) $\rightarrow$ Customer (Feign) $\rightarrow$ BD}\vspace{2.5cm}}}
    \caption{Diagramme de séquence : génération d'une facture}
    \label{fig:seq_invoice}
\end{figure}

\subsection{Architecture technique de la Release 1}

\begin{figure}[H]
    \centering
    \fbox{\parbox{0.85\textwidth}{\centering\vspace{3cm}\textbf{[Diagramme d'architecture microservices]}\\\textit{Frontend Angular $\rightarrow$ API Gateway $\rightarrow$ Eureka $\rightarrow$ Services métier $\rightarrow$ MySQL}\vspace{3cm}}}
    \caption{Architecture technique des microservices}
    \label{fig:architecture}
\end{figure}

\subsection{Interfaces réalisées --- Sprint 3}

\begin{figure}[H]
    \centering
    \fbox{\parbox{0.85\textwidth}{\centering\vspace{2.5cm}\textbf{[Capture d'écran --- Interface de facturation]}\vspace{2.5cm}}}
    \caption{Interface de gestion des factures}
    \label{fig:ui_billing}
\end{figure}

\begin{figure}[H]
    \centering
    \fbox{\parbox{0.85\textwidth}{\centering\vspace{2.5cm}\textbf{[Capture d'écran --- Page de connexion]}\vspace{2.5cm}}}
    \caption{Interface d'authentification}
    \label{fig:ui_login}
\end{figure}

\section*{Conclusion}
\addcontentsline{toc}{section}{Conclusion}
Au terme de cette release, l'ensemble des microservices métier sont opérationnels : gestion des boutiques, des clients, du catalogue, des abonnements et de la facturation. L'authentification JWT et l'API Gateway assurent la sécurité et le routage centralisé. Le chapitre suivant abordera la mise en place de l'infrastructure de déploiement.

\newpage
